%\documentclass[11pt,a4paper,uplatex,dvipdfmx]{ujarticle} 		% for uplatex
\documentclass[11pt,a4j,dvipdfmx]{jarticle} 					% for platex
\input{pieces/form00_header} % pieces
\input{pieces/kakenhi7} % pieces
\input{pieces/form01_header} % pieces
\input{pieces/form02_2023_header} % pieces
\input{pieces/hook3} % pieces
%#Name: gakuhen_b_abst
\input{pieces/form04_jsps_headers} % pieces
\input{pieces/form04_mext_headers_2020} % pieces
\input{pieces/form04_gakuhen_b_abst_header} % pieces
% ===== Global definitions for the Kakenhi form ======================
%　基本情報
%
%------ 研究課題名  -------------------------------------------
\newcommand{\研究課題名}{象の卵}

%----- 研究機関名と研究代表者の氏名-----------------------
\newcommand{\研究機関名}{逢坂大学}
\newcommand{\研究代表者氏名}{湯川秀樹}
\newcommand{\me}{\underline{\underline{H.~Yukawa}}} 
%---- 研究期間の最終年度 ----------------
\newcommand{\研究期間の最終元号年度}{7}  %令和で、半角数字のみ
%========================================

\input{pieces/inst_general} % pieces
\input{pieces/inst_GHb_abst} % pieces
\input{keikaku_defs}
% user07_header
% ===== my favorite packages ====================================
% ここに、自分の使いたいパッケージを宣言して下さい。
\usepackage{wrapfig}
%\usepackage{amssymb}
%\usepackage{mb}
%\DeclareGraphicsRule{.tif}{png}{.png}{`convert #1 `dirname #1`/`basename #1 .tif`.png}
\usepackage{lineno}

% ===== my personal definitions ==================================
% ここに、自分のよく使う記号などを定義して下さい。
\newcommand{\klpionn}{K_L \to \pi^0 \nu \overline{\nu}}
\newcommand{\kppipnn}{K^+ \to \pi^+ \nu \overline{\nu}}

% ----- 業績リスト用 -------------
\newcommand{\paper}[6]{%
	% paper{title}{authors}{journal}{vol}{pages}{year}
	\item ``#1'', #2, #3 {\bf #4}, #5 (#6).			% お好みに合わせて変えてください。
}

\newcommand{\etal}{\textit{et al.\ }}
\newcommand{\ca}[1]{*#1}	% corresponding author;   \ca{\yukawa}  みたいにして使う
\newcommand{\invitedtalk}{招待講演}

\newcommand{\yukawa}{H.~Yukawa}					% no underline
%\newcommand{\yukawa}{\underline{\underline{H.~Yukawa}}}	% with 2 underlines
\newcommand{\tomonaga}{S.~Tomonaga}

\newcommand{\prl}{Phys.\ Rev.\ Lett.\ }		% よく使う雑誌も定義すると楽

% ===== 欄外メモ ==================
\newcommand{\memo}[1]{\marginpar{#1}}
%\renewcommand{\memo}[1]{}	% 全てのメモを表示させないようにするには、行頭の"%"を消す

%%\input{../../pieces/user07_header}	% expand
%\input{../../sample/simple/contents}	% skip
\input{pieces/hook5} % pieces

\begin{document}
\input{pieces/hook7} % pieces
%#Split: 01_plan_abst  
%#PieceName: p01_plan_abst
\input{pieces/p01_plan_abst_00}
\section{領域計画書の概要}
%　　　　＜＜最大　3ページ＞＞

%s08_plan_GHbA
%begin 領域計画書の概要 ＝＝＝＝＝＝＝＝＝＝＝＝＝＝＝＝＝＝＝＝
\JSPSInstructions	% <-- 留意事項。これは消すか、コメントアウトしてください。

\textbf{\\　　　　　＊＊＊　以下は、あくまで例です。真似しないでください。　＊＊＊\\
　　　　　＊＊＊　本文はもちろん、節の切り方や論理の組み方は　　　＊＊＊\\
　　　　　＊＊＊　ご自分の気に入ったスタイルで書いてください。　　＊＊＊}

\subsection{計画研究の名前}
        領域全体の計画書を書くときに用いた\texttt{keikaku\_defs.tex}で定義したコマンド
        (例えば
        \texttt{\textbackslash codeZoo, \textbackslash titleZoo, \textbackslash codeNicknameZoo}など)を用い、
        各計画研究の記号や研究課題名やその略称
        （\codeZoo, \titleZoo, \codeNicknameZoo)
        を表示するのが楽です。
       書き間違えないし、あとで名前が変わっても簡単に直せます。

	領域計画書にはこんなことが書かれてます．．．

    以下の計画研究班を立てる。
    \begin{itemize}
    	\item {\bf \codeTitleZoo} 世界の動物園を巡り、象舍に卵が隠されていないか、探す。
	\item {\bf \codeTitleAfrica} アフリカに行き、空と地上から象の卵を探す。
		アフリカ象は気性が荒いが、サバンナの方がジャングルよりも見通しが効く。
	\item {\bf \codeTitleIndia} インドとタイに行き、ジャングルに隠されている卵を探す。
               ジャングルの場合は空からは探しにくいが、象使いも多く、象の背中に乗って
               象の視点から探索することができる。
               さらに、気だての優しいインド象ならば卵の在処を教えてくれる可能性もある。
    \end{itemize}

	\vspace{1cm}
	\begin{thebibliography}{99}
		\bibitem{teramura} 寺村輝夫、「ぼくは王様 - ぞうのたまごのたまごやき」.
	\end{thebibliography}

	\input{blahblah}
%end 領域計画書の概要 ＝＝＝＝＝＝＝＝＝＝＝＝＝＝＝＝＝＝＝＝
\input{pieces/p01_plan_abst_01}

%#Split: keikaku_GHbA	% keep
\KLInput{keikaku_abst_x00}
\KLInput{keikaku_abst_a01}
\KLInput{keikaku_abst_a02}

% 
% keikaku_abst_z99.tex を複製し、
% z99 をa02などの研究項目番号に改名し、
% そのファイルを上の例のように\KLInput で読み込んでください。
% \KLInputを使うと、読み込むファイルの中で定義されたコマンドはそのファイルの中でのみ
% 有効で、他に影響を及ぼさないので安全です。
%===================================



%#Split: 99_tail
\input{pieces/hook9} % pieces
\end{document}

